% --- Archon physics formalization source begin ---
% archon:physics
% archon:covers PhyXMiniProblems/problem_phyx_mini_0540.lean
% archon:source-report reports/phyx_mini/problem_phyx_mini_0540.source.json

\chapter{Physics problem phyx\_mini\_0540}
\label{ch:PhyXMiniProblems_problem_phyx_mini_0540}

\paragraph{Problem source.}
\#\# Physical scenario

The particle decay \textbackslash{}(\textbackslash{}Sigma\textasciicircum{}+ \textbackslash{}rightarrow \textbackslash{}pi\textasciicircum{}+ + n\textbackslash{}) is observed in a bubble chamber. Figure represents the curved tracks of the particles \textbackslash{}(\textbackslash{}Sigma\textasciicircum{}+\textbackslash{}) and \textbackslash{}(\textbackslash{}pi\textasciicircum{}+\textbackslash{}) and the invisible track of the neutron in the presence of a uniform magnetic field of \textbackslash{}(1.15 \textbackslash{}text\{ T\}\textbackslash{}) directed out of the page. The measured radii of curvature are \textbackslash{}(1.99 \textbackslash{}text\{ m\}\textbackslash{}) for the \textbackslash{}(\textbackslash{}Sigma\textasciicircum{}+\textbackslash{}) particle and \textbackslash{}(0.580 \textbackslash{}text\{ m\}\textbackslash{}) for the \textbackslash{}(\textbackslash{}pi\textasciicircum{}+\textbackslash{}) particle. From this information, we wish to determine the mass of the \textbackslash{}(\textbackslash{}Sigma\textasciicircum{}+\textbackslash{}) particle.

\#\# Question

The angle between the momenta of the \textbackslash{}(\textbackslash{}Sigma\textasciicircum{}+\textbackslash{}) and the \textbackslash{}(\textbackslash{}pi\textasciicircum{}+\textbackslash{}) particles at the moment of decay is \textbackslash{}(\textbackslash{}theta = 64.5\textasciicircum{}\textbackslash{}circ\textbackslash{}). Find the magnitude of the momentum of the neutron.

\#\# Answer choices

- A: A: \textbackslash{}(710 \textbackslash{}text\{ MeV/\}c\textbackslash{})

- B: B: \textbackslash{}(509 \textbackslash{}text\{ MeV/\}c\textbackslash{})

- C: C: \textbackslash{}(533 \textbackslash{}text\{ MeV/\}c\textbackslash{})

- D: D: \textbackslash{}(626 \textbackslash{}text\{ MeV/\}c\textbackslash{})

\#\# Figure caption (auxiliary; use the image as primary evidence)

The image shows a diagram with several vectors and a labeled angle. Here is a detailed description:

1. **Vectors and Lines**:
   - A dashed line labeled "n" extends diagonally from the left.
   - A solid vertical line intersects with the dashed line.
   - Two curved solid lines are emanating from the intersection point. 
     - One is labeled "Σ⁺" with an arrow pointing downward and slightly to the left.
     - The other is labeled "π⁺" with an arrow pointing to the right and slightly upward.

2. **Angle**:
   - An angle is formed between the vertical line and the "π⁺" line. This angle is marked as "θ".

The diagram likely represents a physics or mathematics concept involving vectors and angles.

\#\# Dataset metadata

Particle Physics; Spatial Relation Reasoning, Multi-Formula Reasoning

\paragraph{Recorded answer/context.}
D: D: \textbackslash{}(626 \textbackslash{}text\{ MeV/\}c\textbackslash{})

\paragraph{Figure/image path.}
/root/proposal\_for\_physic/science-mango/phyx\_mini\_run/phyx\_data/test\_image/540.png

\paragraph{Formalization target.}
create a compiling Lean file with sorry bodies at `PhyXMiniProblems/problem\_phyx\_mini\_0540.lean`. The Lean declarations must preserve the physical quantities, dimensions or dimensional roles, figure labels, governing-law hypotheses, and final relation expressed by this problem.
Use Mathlib/Physlib names found through LeanExplore where available. If a physics API is missing, introduce faithful local abstractions rather than scalar placeholder aliases.

\begin{theorem}[Physics formalization target]
\label{thm:physics:phyx_mini_0540:target}
\lean{PhyXMiniProblems.ProblemPhyXMini0540.problem_phyx_mini_0540}
\uses{def:physics:phyx-mini-0540:phyxminiproblems-problemphyxmini0540-sigmaplusdecaysetup, def:physics:phyx-mini-0540:phyxminiproblems-problemphyxmini0540-matchessigmaplusdecayscenario, def:physics:phyx-mini-0540:phyxminiproblems-problemphyxmini0540-matchesproblemreadouts, def:physics:phyx-mini-0540:phyxminiproblems-problemphyxmini0540-matchessuppliedfigure, def:physics:phyx-mini-0540:phyxminiproblems-problemphyxmini0540-usesstandardparticlecharges, def:physics:phyx-mini-0540:phyxminiproblems-problemphyxmini0540-hasphysicaldecayparameters, def:physics:phyx-mini-0540:phyxminiproblems-problemphyxmini0540-isuniformmagneticfieldoutofpage, def:physics:phyx-mini-0540:phyxminiproblems-problemphyxmini0540-satisfieschargedtrackcurvaturelaw, def:physics:phyx-mini-0540:phyxminiproblems-problemphyxmini0540-satisfiesdecaymomentumconservation, def:physics:phyx-mini-0540:phyxminiproblems-problemphyxmini0540-isclosestdisplayedmomentumchoice}
The assigned autoformalize agent should translate this physics problem into Lean declarations in the covered file, with theorem and lemma proof bodies written as `by sorry`.
\end{theorem}
\begin{proof}
This is an autoformalization task, not a proof task. Produce statements that can later be proved by the physics prover without weakening the physical model.
\end{proof}

% --- Archon declaration topology begin ---
\section{Declaration topology}
\begin{definition}[magnetic Field Strength Dimension]
  \label{def:physics:phyx-mini-0540:phyxminiproblems-problemphyxmini0540-magneticfieldstrengthdimension}
  \lean{PhyXMiniProblems.ProblemPhyXMini0540.magneticFieldStrengthDimension}
  The SI dimension M T⁻¹ C⁻¹ of magnetic flux density
\end{definition}
\begin{proof}
  This is the declared model definition of magnetic Field Strength Dimension.
\end{proof}
\begin{definition}[Length Quantity]
  \label{def:physics:phyx-mini-0540:phyxminiproblems-problemphyxmini0540-lengthquantity}
  \lean{PhyXMiniProblems.ProblemPhyXMini0540.LengthQuantity}
  A nonnegative, unit-independent physical length
\end{definition}
\begin{proof}
  This is the declared model definition of Length Quantity.
\end{proof}
\begin{definition}[Mass Quantity]
  \label{def:physics:phyx-mini-0540:phyxminiproblems-problemphyxmini0540-massquantity}
  \lean{PhyXMiniProblems.ProblemPhyXMini0540.MassQuantity}
  A nonnegative, unit-independent rest mass
\end{definition}
\begin{proof}
  This is the declared model definition of Mass Quantity.
\end{proof}
\begin{definition}[Charge Magnitude Quantity]
  \label{def:physics:phyx-mini-0540:phyxminiproblems-problemphyxmini0540-chargemagnitudequantity}
  \lean{PhyXMiniProblems.ProblemPhyXMini0540.ChargeMagnitudeQuantity}
  A nonnegative electric-charge magnitude
\end{definition}
\begin{proof}
  This is the declared model definition of Charge Magnitude Quantity.
\end{proof}
\begin{definition}[Magnetic Field Strength Quantity]
  \label{def:physics:phyx-mini-0540:phyxminiproblems-problemphyxmini0540-magneticfieldstrengthquantity}
  \lean{PhyXMiniProblems.ProblemPhyXMini0540.MagneticFieldStrengthQuantity}
  \uses{def:physics:phyx-mini-0540:phyxminiproblems-problemphyxmini0540-magneticfieldstrengthdimension}
  A nonnegative magnetic-field-strength magnitude
\end{definition}
\begin{proof}
  This is the declared model definition of Magnetic Field Strength Quantity.
\end{proof}
\begin{definition}[Planar Momentum Quantity]
  \label{def:physics:phyx-mini-0540:phyxminiproblems-problemphyxmini0540-planarmomentumquantity}
  \lean{PhyXMiniProblems.ProblemPhyXMini0540.PlanarMomentumQuantity}
  A unit-independent physical two-momentum. Physlib's Momentum 2 has the correct M L T⁻¹ dimension and stores its two components in the chosen unit system
\end{definition}
\begin{proof}
  This is the declared model definition of Planar Momentum Quantity.
\end{proof}
\begin{definition}[length In Meters]
  \label{def:physics:phyx-mini-0540:phyxminiproblems-problemphyxmini0540-lengthinmeters}
  \lean{PhyXMiniProblems.ProblemPhyXMini0540.lengthInMeters}
  \uses{def:physics:phyx-mini-0540:phyxminiproblems-problemphyxmini0540-lengthquantity}
  Read a physical length in coherent SI metres
\end{definition}
\begin{proof}
  This is the declared model definition of length In Meters.
\end{proof}
\begin{definition}[charge In Coulombs]
  \label{def:physics:phyx-mini-0540:phyxminiproblems-problemphyxmini0540-chargeincoulombs}
  \lean{PhyXMiniProblems.ProblemPhyXMini0540.chargeInCoulombs}
  \uses{def:physics:phyx-mini-0540:phyxminiproblems-problemphyxmini0540-chargemagnitudequantity}
  Read a charge magnitude in coherent SI coulombs
\end{definition}
\begin{proof}
  This is the declared model definition of charge In Coulombs.
\end{proof}
\begin{definition}[magnetic Field Strength In Teslas]
  \label{def:physics:phyx-mini-0540:phyxminiproblems-problemphyxmini0540-magneticfieldstrengthinteslas}
  \lean{PhyXMiniProblems.ProblemPhyXMini0540.magneticFieldStrengthInTeslas}
  \uses{def:physics:phyx-mini-0540:phyxminiproblems-problemphyxmini0540-magneticfieldstrengthquantity}
  Read magnetic flux density in coherent SI teslas
\end{definition}
\begin{proof}
  This is the declared model definition of magnetic Field Strength In Teslas.
\end{proof}
\begin{definition}[momentum Vector Readout]
  \label{def:physics:phyx-mini-0540:phyxminiproblems-problemphyxmini0540-momentumvectorreadout}
  \lean{PhyXMiniProblems.ProblemPhyXMini0540.momentumVectorReadout}
  \uses{def:physics:phyx-mini-0540:phyxminiproblems-problemphyxmini0540-planarmomentumquantity}
  Read a physical two-momentum in arbitrary coherent units and equip its two components with Mathlib's Euclidean norm and inner product
\end{definition}
\begin{proof}
  This is the declared model definition of momentum Vector Readout.
\end{proof}
\begin{definition}[momentum Vector In SI]
  \label{def:physics:phyx-mini-0540:phyxminiproblems-problemphyxmini0540-momentumvectorinsi}
  \lean{PhyXMiniProblems.ProblemPhyXMini0540.momentumVectorInSI}
  \uses{def:physics:phyx-mini-0540:phyxminiproblems-problemphyxmini0540-planarmomentumquantity, def:physics:phyx-mini-0540:phyxminiproblems-problemphyxmini0540-momentumvectorreadout}
  Planar momentum-vector readout in kilogram-metres per second
\end{definition}
\begin{proof}
  This is the declared model definition of momentum Vector In SI.
\end{proof}
\begin{definition}[momentum Magnitude In Kilogram Meters Per Second]
  \label{def:physics:phyx-mini-0540:phyxminiproblems-problemphyxmini0540-momentummagnitudeinkilogrammeterspersecond}
  \lean{PhyXMiniProblems.ProblemPhyXMini0540.momentumMagnitudeInKilogramMetersPerSecond}
  \uses{def:physics:phyx-mini-0540:phyxminiproblems-problemphyxmini0540-planarmomentumquantity, def:physics:phyx-mini-0540:phyxminiproblems-problemphyxmini0540-momentumvectorinsi}
  Momentum magnitude in kilogram-metres per second
\end{definition}
\begin{proof}
  This is the declared model definition of momentum Magnitude In Kilogram Meters Per Second.
\end{proof}
\begin{definition}[speed Of Light In Meters Per Second]
  \label{def:physics:phyx-mini-0540:phyxminiproblems-problemphyxmini0540-speedoflightinmeterspersecond}
  \lean{PhyXMiniProblems.ProblemPhyXMini0540.speedOfLightInMetersPerSecond}
  The SI metre-per-second readout of Physlib's exact speed of light
\end{definition}
\begin{proof}
  This is the declared model definition of speed Of Light In Meters Per Second.
\end{proof}
\begin{definition}[electron Volt In Joules]
  \label{def:physics:phyx-mini-0540:phyxminiproblems-problemphyxmini0540-electronvoltinjoules}
  \lean{PhyXMiniProblems.ProblemPhyXMini0540.electronVoltInJoules}
  The SI joule readout of Physlib's exact electron volt
\end{definition}
\begin{proof}
  This is the declared model definition of electron Volt In Joules.
\end{proof}
\begin{definition}[momentum Magnitude In Me VPer C]
  \label{def:physics:phyx-mini-0540:phyxminiproblems-problemphyxmini0540-momentummagnitudeinmevperc}
  \lean{PhyXMiniProblems.ProblemPhyXMini0540.momentumMagnitudeInMeVPerC}
  \uses{def:physics:phyx-mini-0540:phyxminiproblems-problemphyxmini0540-planarmomentumquantity, def:physics:phyx-mini-0540:phyxminiproblems-problemphyxmini0540-momentummagnitudeinkilogrammeterspersecond, def:physics:phyx-mini-0540:phyxminiproblems-problemphyxmini0540-speedoflightinmeterspersecond, def:physics:phyx-mini-0540:phyxminiproblems-problemphyxmini0540-electronvoltinjoules}
  Convert an SI momentum magnitude to MeV/c, using p[MeV/c] = p[kg m/s] * c / (10\textasciicircum{}6 eV[J])
\end{definition}
\begin{proof}
  This is the declared model definition of momentum Magnitude In Me VPer C.
\end{proof}
\begin{definition}[Particle Species]
  \label{def:physics:phyx-mini-0540:phyxminiproblems-problemphyxmini0540-particlespecies}
  \lean{PhyXMiniProblems.ProblemPhyXMini0540.ParticleSpecies}
  The three particle species participating in the observed decay
\end{definition}
\begin{proof}
  This is the declared model definition of Particle Species.
\end{proof}
\begin{definition}[Charged Track]
  \label{def:physics:phyx-mini-0540:phyxminiproblems-problemphyxmini0540-chargedtrack}
  \lean{PhyXMiniProblems.ProblemPhyXMini0540.ChargedTrack}
  Charged tracks whose curvature radii are measured in the chamber
\end{definition}
\begin{proof}
  This is the declared model definition of Charged Track.
\end{proof}
\begin{definition}[particle]
  \label{def:physics:phyx-mini-0540:phyxminiproblems-problemphyxmini0540-chargedtrack-particle}
  \lean{PhyXMiniProblems.ProblemPhyXMini0540.ChargedTrack.particle}
  \uses{def:physics:phyx-mini-0540:phyxminiproblems-problemphyxmini0540-particlespecies, def:physics:phyx-mini-0540:phyxminiproblems-problemphyxmini0540-chargedtrack}
  The particle represented by each charged track
\end{definition}
\begin{proof}
  This is the declared model definition of particle.
\end{proof}
\begin{definition}[Charge Sign]
  \label{def:physics:phyx-mini-0540:phyxminiproblems-problemphyxmini0540-chargesign}
  \lean{PhyXMiniProblems.ProblemPhyXMini0540.ChargeSign}
  The sign class of a particle's electric charge
\end{definition}
\begin{proof}
  This is the declared model definition of Charge Sign.
\end{proof}
\begin{definition}[Observation Medium]
  \label{def:physics:phyx-mini-0540:phyxminiproblems-problemphyxmini0540-observationmedium}
  \lean{PhyXMiniProblems.ProblemPhyXMini0540.ObservationMedium}
  The apparatus medium in which the tracks are observed
\end{definition}
\begin{proof}
  This is the declared model definition of Observation Medium.
\end{proof}
\begin{definition}[Page Normal Direction]
  \label{def:physics:phyx-mini-0540:phyxminiproblems-problemphyxmini0540-pagenormaldirection}
  \lean{PhyXMiniProblems.ProblemPhyXMini0540.PageNormalDirection}
  Directions normal to the two-dimensional page
\end{definition}
\begin{proof}
  This is the declared model definition of Page Normal Direction.
\end{proof}
\begin{definition}[Figure Label]
  \label{def:physics:phyx-mini-0540:phyxminiproblems-problemphyxmini0540-figurelabel}
  \lean{PhyXMiniProblems.ProblemPhyXMini0540.FigureLabel}
  Text and symbols visibly printed in the supplied raster
\end{definition}
\begin{proof}
  This is the declared model definition of Figure Label.
\end{proof}
\begin{definition}[Figure Element]
  \label{def:physics:phyx-mini-0540:phyxminiproblems-problemphyxmini0540-figureelement}
  \lean{PhyXMiniProblems.ProblemPhyXMini0540.FigureElement}
  Geometric elements visible in the supplied raster
\end{definition}
\begin{proof}
  This is the declared model definition of Figure Element.
\end{proof}
\begin{definition}[Figure Stroke Style]
  \label{def:physics:phyx-mini-0540:phyxminiproblems-problemphyxmini0540-figurestrokestyle}
  \lean{PhyXMiniProblems.ProblemPhyXMini0540.FigureStrokeStyle}
  Stroke styles used to distinguish tracks and inferred rays
\end{definition}
\begin{proof}
  This is the declared model definition of Figure Stroke Style.
\end{proof}
\begin{definition}[Arrow Flow At Vertex]
  \label{def:physics:phyx-mini-0540:phyxminiproblems-problemphyxmini0540-arrowflowatvertex}
  \lean{PhyXMiniProblems.ProblemPhyXMini0540.ArrowFlowAtVertex}
  Whether a momentum arrow enters or leaves the decay vertex
\end{definition}
\begin{proof}
  This is the declared model definition of Arrow Flow At Vertex.
\end{proof}
\begin{definition}[Bubble Chamber Decay Figure]
  \label{def:physics:phyx-mini-0540:phyxminiproblems-problemphyxmini0540-bubblechamberdecayfigure}
  \lean{PhyXMiniProblems.ProblemPhyXMini0540.BubbleChamberDecayFigure}
  \uses{def:physics:phyx-mini-0540:phyxminiproblems-problemphyxmini0540-particlespecies, def:physics:phyx-mini-0540:phyxminiproblems-problemphyxmini0540-figurelabel, def:physics:phyx-mini-0540:phyxminiproblems-problemphyxmini0540-figureelement, def:physics:phyx-mini-0540:phyxminiproblems-problemphyxmini0540-figurestrokestyle, def:physics:phyx-mini-0540:phyxminiproblems-problemphyxmini0540-arrowflowatvertex}
  Presentation-level data transcribed from image 540. The dashed neutron line is a reconstructed momentum ray, not an ionization track in the chamber
\end{definition}
\begin{proof}
  This is the declared model definition of Bubble Chamber Decay Figure.
\end{proof}
\begin{definition}[out Of Page Direction]
  \label{def:physics:phyx-mini-0540:phyxminiproblems-problemphyxmini0540-outofpagedirection}
  \lean{PhyXMiniProblems.ProblemPhyXMini0540.outOfPageDirection}
  The unit vector normal to the page in its positive, outward direction
\end{definition}
\begin{proof}
  This is the declared model definition of out Of Page Direction.
\end{proof}
\begin{definition}[Sigma Plus Decay Setup]
  \label{def:physics:phyx-mini-0540:phyxminiproblems-problemphyxmini0540-sigmaplusdecaysetup}
  \lean{PhyXMiniProblems.ProblemPhyXMini0540.SigmaPlusDecaySetup}
  \uses{def:physics:phyx-mini-0540:phyxminiproblems-problemphyxmini0540-lengthquantity, def:physics:phyx-mini-0540:phyxminiproblems-problemphyxmini0540-massquantity, def:physics:phyx-mini-0540:phyxminiproblems-problemphyxmini0540-chargemagnitudequantity, def:physics:phyx-mini-0540:phyxminiproblems-problemphyxmini0540-magneticfieldstrengthquantity, def:physics:phyx-mini-0540:phyxminiproblems-problemphyxmini0540-planarmomentumquantity, def:physics:phyx-mini-0540:phyxminiproblems-problemphyxmini0540-particlespecies, def:physics:phyx-mini-0540:phyxminiproblems-problemphyxmini0540-chargedtrack, def:physics:phyx-mini-0540:phyxminiproblems-problemphyxmini0540-chargesign, def:physics:phyx-mini-0540:phyxminiproblems-problemphyxmini0540-observationmedium, def:physics:phyx-mini-0540:phyxminiproblems-problemphyxmini0540-pagenormaldirection, def:physics:phyx-mini-0540:phyxminiproblems-problemphyxmini0540-bubblechamberdecayfigure}
  All independent physical quantities in the decay setup. Rest masses are retained because the surrounding scenario aims eventually to determine the Sigma+ mass, although the present subquestion asks only for neutron momentum. No numerical neutron momentum is stored here
\end{definition}
\begin{proof}
  This is the declared model definition of Sigma Plus Decay Setup.
\end{proof}
\begin{definition}[Matches Sigma Plus Decay Scenario]
  \label{def:physics:phyx-mini-0540:phyxminiproblems-problemphyxmini0540-matchessigmaplusdecayscenario}
  \lean{PhyXMiniProblems.ProblemPhyXMini0540.MatchesSigmaPlusDecayScenario}
  \uses{def:physics:phyx-mini-0540:phyxminiproblems-problemphyxmini0540-sigmaplusdecaysetup}
  Categorical content of the physical scenario. The two charged particles are positive, the neutron is neutral and invisible in the chamber, and all momenta lie in the page perpendicular to the field
\end{definition}
\begin{proof}
  This is the declared model definition of Matches Sigma Plus Decay Scenario.
\end{proof}
\begin{definition}[Matches Problem Readouts]
  \label{def:physics:phyx-mini-0540:phyxminiproblems-problemphyxmini0540-matchesproblemreadouts}
  \lean{PhyXMiniProblems.ProblemPhyXMini0540.MatchesProblemReadouts}
  \uses{def:physics:phyx-mini-0540:phyxminiproblems-problemphyxmini0540-lengthinmeters, def:physics:phyx-mini-0540:phyxminiproblems-problemphyxmini0540-magneticfieldstrengthinteslas, def:physics:phyx-mini-0540:phyxminiproblems-problemphyxmini0540-momentumvectorinsi, def:physics:phyx-mini-0540:phyxminiproblems-problemphyxmini0540-sigmaplusdecaysetup}
  Numerical measurements stated in the problem. The field is 1.15 T, the curvature radii are 1.99 m and 0.580 m, and the momentum angle is 64.5 degrees = 43*pi/120 radians. No neutron momentum value occurs here
\end{definition}
\begin{proof}
  This is the declared model definition of Matches Problem Readouts.
\end{proof}
\begin{definition}[Matches Supplied Figure]
  \label{def:physics:phyx-mini-0540:phyxminiproblems-problemphyxmini0540-matchessuppliedfigure}
  \lean{PhyXMiniProblems.ProblemPhyXMini0540.MatchesSuppliedFigure}
  \uses{def:physics:phyx-mini-0540:phyxminiproblems-problemphyxmini0540-sigmaplusdecaysetup}
  Qualitative geometry read directly from the primary image. In particular, the incoming Sigma+ arrow terminates at the decay vertex, while the pion and reconstructed neutron arrows leave it
\end{definition}
\begin{proof}
  This is the declared model definition of Matches Supplied Figure.
\end{proof}
\begin{definition}[Uses Standard Particle Charges]
  \label{def:physics:phyx-mini-0540:phyxminiproblems-problemphyxmini0540-usesstandardparticlecharges}
  \lean{PhyXMiniProblems.ProblemPhyXMini0540.UsesStandardParticleCharges}
  \uses{def:physics:phyx-mini-0540:phyxminiproblems-problemphyxmini0540-chargeincoulombs, def:physics:phyx-mini-0540:phyxminiproblems-problemphyxmini0540-sigmaplusdecaysetup}
  Standard elementary-charge calibration. Both positive charged particles have magnitude e, while the neutral neutron has zero charge magnitude
\end{definition}
\begin{proof}
  This is the declared model definition of Uses Standard Particle Charges.
\end{proof}
\begin{definition}[Has Physical Decay Parameters]
  \label{def:physics:phyx-mini-0540:phyxminiproblems-problemphyxmini0540-hasphysicaldecayparameters}
  \lean{PhyXMiniProblems.ProblemPhyXMini0540.HasPhysicalDecayParameters}
  \uses{def:physics:phyx-mini-0540:phyxminiproblems-problemphyxmini0540-lengthinmeters, def:physics:phyx-mini-0540:phyxminiproblems-problemphyxmini0540-chargeincoulombs, def:physics:phyx-mini-0540:phyxminiproblems-problemphyxmini0540-magneticfieldstrengthinteslas, def:physics:phyx-mini-0540:phyxminiproblems-problemphyxmini0540-momentumvectorinsi, def:physics:phyx-mini-0540:phyxminiproblems-problemphyxmini0540-sigmaplusdecaysetup}
  Positivity and nondegeneracy conditions for the observed event
\end{definition}
\begin{proof}
  This is the declared model definition of Has Physical Decay Parameters.
\end{proof}
\begin{definition}[Is Uniform Magnetic Field Out Of Page]
  \label{def:physics:phyx-mini-0540:phyxminiproblems-problemphyxmini0540-isuniformmagneticfieldoutofpage}
  \lean{PhyXMiniProblems.ProblemPhyXMini0540.IsUniformMagneticFieldOutOfPage}
  \uses{def:physics:phyx-mini-0540:phyxminiproblems-problemphyxmini0540-magneticfieldstrengthinteslas, def:physics:phyx-mini-0540:phyxminiproblems-problemphyxmini0540-outofpagedirection, def:physics:phyx-mini-0540:phyxminiproblems-problemphyxmini0540-sigmaplusdecaysetup}
  Physlib's electromagnetic field uses real vector coordinates. This premise states that those coordinates are tesla readouts and that the field is the uniform vector B z-hat at every spacetime point
\end{definition}
\begin{proof}
  This is the declared model definition of Is Uniform Magnetic Field Out Of Page.
\end{proof}
\begin{definition}[Satisfies Charged Track Curvature Law]
  \label{def:physics:phyx-mini-0540:phyxminiproblems-problemphyxmini0540-satisfieschargedtrackcurvaturelaw}
  \lean{PhyXMiniProblems.ProblemPhyXMini0540.SatisfiesChargedTrackCurvatureLaw}
  \uses{def:physics:phyx-mini-0540:phyxminiproblems-problemphyxmini0540-lengthinmeters, def:physics:phyx-mini-0540:phyxminiproblems-problemphyxmini0540-chargeincoulombs, def:physics:phyx-mini-0540:phyxminiproblems-problemphyxmini0540-magneticfieldstrengthinteslas, def:physics:phyx-mini-0540:phyxminiproblems-problemphyxmini0540-momentummagnitudeinkilogrammeterspersecond, def:physics:phyx-mini-0540:phyxminiproblems-problemphyxmini0540-sigmaplusdecaysetup}
  For each positive charged track perpendicular to a uniform field, the relativistic momentum magnitude obeys p = q B R. This is the governing magnetic-curvature law, not a specialized neutron-momentum formula
\end{definition}
\begin{proof}
  This is the declared model definition of Satisfies Charged Track Curvature Law.
\end{proof}
\begin{definition}[Satisfies Decay Momentum Conservation]
  \label{def:physics:phyx-mini-0540:phyxminiproblems-problemphyxmini0540-satisfiesdecaymomentumconservation}
  \lean{PhyXMiniProblems.ProblemPhyXMini0540.SatisfiesDecayMomentumConservation}
  \uses{def:physics:phyx-mini-0540:phyxminiproblems-problemphyxmini0540-momentumvectorreadout, def:physics:phyx-mini-0540:phyxminiproblems-problemphyxmini0540-sigmaplusdecaysetup}
  Three-momentum conservation at the two-body decay vertex, restricted to the two-dimensional chamber plane. It is required in every coherent unit system and contains no numerical neutron answer
\end{definition}
\begin{proof}
  This is the declared model definition of Satisfies Decay Momentum Conservation.
\end{proof}
\begin{lemma}[neutron Momentum Cosine Relation]
  \label{lem:physics:phyx-mini-0540:phyxminiproblems-problemphyxmini0540-neutronmomentumcosinerelation}
  \lean{PhyXMiniProblems.ProblemPhyXMini0540.neutronMomentumCosineRelation}
  \uses{def:physics:phyx-mini-0540:phyxminiproblems-problemphyxmini0540-momentummagnitudeinkilogrammeterspersecond, def:physics:phyx-mini-0540:phyxminiproblems-problemphyxmini0540-sigmaplusdecaysetup, def:physics:phyx-mini-0540:phyxminiproblems-problemphyxmini0540-matchesproblemreadouts, def:physics:phyx-mini-0540:phyxminiproblems-problemphyxmini0540-hasphysicaldecayparameters, def:physics:phyx-mini-0540:phyxminiproblems-problemphyxmini0540-satisfiesdecaymomentumconservation}
  Momentum conservation and the measured parent-pion angle imply the cosine rule for the neutron momentum. This is a derived intermediate conclusion, not a field of any premise structure
\end{lemma}
\begin{proof}
  The conclusion follows from the governing relations and figure readouts named in the statement of neutron Momentum Cosine Relation.
\end{proof}
\begin{definition}[Answer Choice]
  \label{def:physics:phyx-mini-0540:phyxminiproblems-problemphyxmini0540-answerchoice}
  \lean{PhyXMiniProblems.ProblemPhyXMini0540.AnswerChoice}
  Labels of the four momentum values displayed with the question
\end{definition}
\begin{proof}
  This is the declared model definition of Answer Choice.
\end{proof}
\begin{definition}[displayed Momentum In Me VPer C]
  \label{def:physics:phyx-mini-0540:phyxminiproblems-problemphyxmini0540-displayedmomentuminmevperc}
  \lean{PhyXMiniProblems.ProblemPhyXMini0540.displayedMomentumInMeVPerC}
  \uses{def:physics:phyx-mini-0540:phyxminiproblems-problemphyxmini0540-answerchoice}
  The displayed neutron-momentum candidates, as numerical MeV/c values
\end{definition}
\begin{proof}
  This is the declared model definition of displayed Momentum In Me VPer C.
\end{proof}
\begin{definition}[Is Closest Displayed Momentum Choice]
  \label{def:physics:phyx-mini-0540:phyxminiproblems-problemphyxmini0540-isclosestdisplayedmomentumchoice}
  \lean{PhyXMiniProblems.ProblemPhyXMini0540.IsClosestDisplayedMomentumChoice}
  \uses{def:physics:phyx-mini-0540:phyxminiproblems-problemphyxmini0540-momentummagnitudeinmevperc, def:physics:phyx-mini-0540:phyxminiproblems-problemphyxmini0540-sigmaplusdecaysetup, def:physics:phyx-mini-0540:phyxminiproblems-problemphyxmini0540-answerchoice, def:physics:phyx-mini-0540:phyxminiproblems-problemphyxmini0540-displayedmomentuminmevperc}
  A displayed choice is correct when its MeV/c readout is at least as close to the physically derived neutron momentum as every other displayed choice. This comparison is appropriate for the rounded measured inputs and avoids asserting that the underlying physical magnitude equals a whole number exactly
\end{definition}
\begin{proof}
  This is the declared model definition of Is Closest Displayed Momentum Choice.
\end{proof}
% --- Archon declaration topology end ---
% --- Archon physics formalization source end ---
